% Section: The analysis framework
\section{The analysis framework}
\label{sec:analysis}

The analysis libraries (\code{src/analysis}) build on the core to process events and fill histograms. The
\code{Particles} library holds what all analyses share (events, particles, particle databases, filters,
efficiencies, smearers, the base classes of the analysis tasks and of their histogram folders), and each other library
holds one family of observables (Section~\ref{sec:observables}). This section describes how an analysis is run.

% ---------------------------------------------------------------------------------------------------------------
\subsection{Tasks and event processing}
\label{sec:tasks}

\paragraph{Jobs.}
A job is one run of the \code{CAP} program on one task tree, read from a workflow file (Section~\ref{sec:xml}):
{\footnotesize
\begin{verbatim}
CAP <task> <configuration directory>/
    <workflow file> <output directory>/
\end{verbatim}}
The program creates the tree of tasks named on the command line, configures, initialises and executes it, and
finalises it (Section~\ref{sec:objects}). Histograms are read and written once per job, not by each task: if the
main task names an input file (\code{JOB:INPUT\_FILE}), the histograms that its tasks import are read from it after
initialisation, and the histograms of all the tasks are written at the end to one output file
(\code{JOB:OUTPUT\_FILE}), in a folder per task (Section~\ref{sec:folders-analysis}).

\paragraph{Analysis tasks and event processors.}
The tasks of an analysis manage lists of the objects they work with, each configured with the keys of its type. An
object of a list is either owned by the task, which creates or imports it and may export it to a file, or shared: it
is then taken by name from a store common to all the tasks of the job, so that, for instance, all the analyses of a
job use the same filters (Section~\ref{sec:filters}). The objects are split between two base classes, so that a task
holds only what it uses. \code{AnalysisTask} holds what every task of an analysis needs: the event filters and the
particle filters, their names, and the naming of the combinations of filters by which the folders of histograms are
found in the files of the jobs (Section~\ref{sec:filters}). \code{EventProcessor}, derived from it, adds what is needed
to process events: the events, the event and particle efficiencies, the event and particle smearers, the particle
databases, and the filtering of each event and of its particles. Event generators and readers, decayers and the
analyzers are event processors; the calculators of the later phases of an analysis (Section~\ref{sec:phases}), which
only read and complete histograms, derive from \code{AnalysisTask} alone, and their configurations accordingly
contain only their filters and histogram settings.

\paragraph{Event loop.}
\code{EventIterator} is the task that runs the event loop: for each of the events requested, it executes its subtasks
in order, typically one or several generators followed by the analyzers, and reports its progress periodically. It
checks when it is initialised that it uses the events that its subtasks make, so that a generator and the analyzers
that follow it work on the same event. The iterator resets the event before each pass, and the generators add their
particles to it, so that several generators can fill one event: a hard process embedded in a soft background, for
example. \code{EventMerger} goes further and fills one event with the particles of several events of each of its input
streams, to simulate pile-up. An \code{Event} holds its particles, created by a particle factory and reused from
event to event to avoid allocations, the two colliding projectiles, and the selections of its particles made by the
particle filters of the tasks, which are shared by all of them (Section~\ref{sec:filters}).

% ---------------------------------------------------------------------------------------------------------------
\subsection{Event and particle filters}
\label{sec:filters}

\paragraph{Filters.}
A filter is a list of conditions, all of which must be satisfied; a filter without conditions accepts everything. A
condition sets a range, two ranges, or a value of one quantity. Event filters select events by a global quantity of the
particles accepted by a particle filter, named by the condition: their multiplicity, energy, net charge, net
strangeness, net baryon number, $\pT$ sum or spherocity; they thereby define event classes such as classes of
centrality or of multiplicity in a given acceptance. Particle filters select particles by their kinematics ($\pT$, the components of the
momentum, the energy, the pseudorapidity, the rapidity and the azimuth), by the properties of their type (PDG code,
charge, strangeness, baryon number, charm and beauty), and by their state. A filter is configured by keys
\code{<owner>:PARTICLE\_FILTER\_<index>:\dots}, for instance
{\footnotesize
\begin{verbatim}
...:CONDITION:N = 2
...:CONDITION_0:TYPE = DOUBLE_RANGE
...:CONDITION_0:SUBTYPE = PT
...:CONDITION_0:MINIMUM = 0.2
...:CONDITION_0:MAXIMUM = 2.0
\end{verbatim}}
\noindent and a condition of an event filter names the particle filter whose selection it tests:
{\footnotesize
\begin{verbatim}
...:CONDITION_0:SUBTYPE = MULT
...:CONDITION_0:SELECTION = Charged
\end{verbatim}}
\noindent Filters are usually not written by hand: they are defined in the filter editor of RunCAP.py (Section~\ref{sec:run}),
stored in the \code{RUNCAP} element of the workflow file, and created at the start of a job by creator tasks
(\code{EventFilterCreator}, \code{ParticleFilterCreator}), which put them in the store from which the other tasks take
them by name. Creators also exist for the efficiencies and the smearers.

\paragraph{Combinations of filters.}
An analysis runs for every combination of its filters. For each event, an event processor records which event filters
accept the event, and, for each particle filter, the list of the particles it accepts and the moments
$\sum_i p_{\mathrm{T},i}^k$ of their transverse momenta, which the analyses of transverse-momentum correlations use.
A single-particle analysis with $n_E$ event filters and $n_P$ particle filters then fills $n_E n_P$ folders of
histograms, one per pair (event filter, particle filter), and a pair analysis fills $n_E n_P^2$ folders of pair
histograms, one per ordered pair of particle filters for each event filter, in addition to its $n_E n_P$ folders of
single-particle histograms. The folder of event filter $e$ and particle filters $p_1, \dots, p_r$ has index
\begin{equation}
  k = e\, n_P^{\,r} + p_1\, n_P^{\,r-1} + \dots + p_r ,
\end{equation}
and is named after the filters, \code{ALL\_PiP} for the positive pions of all events, \code{ALL\_PiP\_PiM} for the pairs of a
positive and a negative pion. When the two filters of a pair are the same, a particle is not paired with itself. The
histograms of each folder are normalised to the number of events accepted by its event filter at the end of the event
loop. One job thus computes, in one pass over the events, the observables of all the species and all the event classes of
an analysis.

\paragraph{Shared selections.}
Several analyses of a job often use the same filters. The selection of a particle filter, the particles it accepts in
an event with their $\pT$ moments and global quantities (\code{ParticleSelection}), and the result of an event filter
are therefore made once per event and kept by the event, which all the tasks share: the first task that needs them
makes them, the others reuse them, and the cost of the filtering no longer grows with the number of analyses. Each
filter receives, when the tasks are initialised, a slot, the index of its selection in the events; a task maps its own
filters to their slots, so that reaching a selection takes two array lookups. Filters are identified by their object:
those of the store, shared by the tasks, share a slot, and a filter made by a task for itself has a slot of its own, so
that analyses can mix shared and private filters freely. The selections are made again whenever the particles of the
event change: a new event, particles added (the products of decays), or particles changed in place (by the tasks that
rotate or shift the events). Since an event filter computes the quantity it tests from the selection it names, it no
longer depends on another task, such as \code{GlobalAnalyzer}, having run before it.

% ---------------------------------------------------------------------------------------------------------------
\subsection{Efficiencies and Signal Smearing}
\label{sec:efficiency}

\paragraph{Efficiencies.}
A particle efficiency (\code{ParticleEfficiency}) is a histogram of the detection efficiency as a function of
pseudorapidity and $\pT$. It is used in two ways: to simulate detection losses, by accepting each particle with the
probability given by the efficiency, and to correct for them, by weighting each particle by the inverse of the
efficiency. The single-particle analyzer weights the particles it counts in this way, with either one efficiency for all
particle filters or one per filter. An event efficiency (\code{EventEfficiency}) has the same role for events, as a
function of their multiplicity.
\todo{the pair analyzers do not weight pairs by the product of the particle efficiencies, and \code{EventEfficiency}
always accepts events (efficiency 1): implement both, or describe them as planned.}

\paragraph{Smearing.}
A particle smearer (\code{ParticleSmearer}) simulates the resolution of a detector: it draws the $\pT$, the
pseudorapidity and the azimuth of a particle from Gaussian distributions centred on the true values plus a bias, with
a width; the bias and the width of each variable are histograms as functions of pseudorapidity and $\pT$. The task
\code{MeasurementSimulator} uses the smearers to produce, from each generated event, a detector-level event: the
particles accepted by the filters are smeared, separately for each combination of filters, and added to a second
event. The same analysis can then run on both events: \code{PerformanceAnalyzer} fills histograms comparing them, and
\code{ClosureCalculator} compares the results of an analysis at the two levels, histogram by histogram, as differences
or ratios (closure tests, Section~\ref{sec:performance}). These tools are meant for fast studies of the performance of
a measurement: of the effect of an acceptance, an efficiency or a resolution on an observable, before a full simulation
of the detector.

% ---------------------------------------------------------------------------------------------------------------
\subsection{Three phases: analyzers, derived calculators, balance functions}
\label{sec:phases}

Each analysis is split into tasks that run in separate jobs, one after the other, each reading the output file of the
previous one.

\paragraph{Analyzers.}
The analyzer of an analysis (\code{ParticleSingleAnalyzer}, \code{ParticlePairAnalyzer}, \code{PtPtAnalyzer},
\code{NuDynAnalyzer}, \dots) runs in the event loop and fills the base histograms event by event: the single-particle
densities, the pair densities, and the sums weighted by the transverse momenta from which the correlators are computed.
It does nothing else, so that the event loop, where the time of a job is spent, is as short as possible, and so that
the base histograms of several jobs can be added before anything is computed from them.

\paragraph{Derived calculators.}
The derived calculator of an analysis (\code{ParticleSingleDerivedCalculator}, \code{ParticlePairDerivedCalculator},
\code{ParticlePair3DDerivedCalculator}, \code{PtPtDerivedCalculator}, \dots) runs in a later job, as a task of its
own without event loop: it reads the base histograms from the output file of the analyzers and computes the
derived histograms: the normalised densities, the products of single-particle densities, the cumulants $C_2$ and
$R_2$, the transverse-momentum correlators, and the projections on $(\Delta\eta,\Delta\varphi)$ and
$(\Delta y,\Delta\varphi)$ (Section~\ref{sec:observables}). For a pair analysis, the derived histograms of each pair folder $(e, i, j)$ are computed from it and from the single-particle folders $(e, i)$ and $(e, j)$. The calculator uses
the same filters and histogram settings as its analyzer, and reads and writes the folders of the analyzer, named by
the key \code{SOURCE\_TASK}: it completes them rather than writing folders of its own, so that the base and derived
histograms of a combination of filters stay together. The base and derived histograms of a combination of filters are
held by one folder class, but the derived histograms are created only by the calculator, so that the analyzer neither
allocates nor writes them. When a folder is read from a file, its derived histograms are read with it if the file has
them, and are otherwise left to be computed: the same folders thus read the output of an analyzer and that of a
calculator. The derived histograms can be recomputed, with other settings or a correction, without running the events
again.

\paragraph{Balance functions.}
A third phase computes the balance functions of a pair analysis (\code{ParticlePairBFCalculator},
\code{ParticlePair3DBFCalculator}, Section~\ref{sec:bf}). Like the derived calculator, it reads the folders of the
analyzer named by \code{SOURCE\_TASK}, with the same filters and histogram settings, and computes their derived
histograms; its input can thus be the output of the analyzer or of the derived calculator. A species is a particle
filter $a$ and the filter $\bar a$ of its antiparticle, named by the key \code{SPECIES}
(e.g.\ \code{PiP/PiM, KP/KM}). For each event filter $e$ and ordered pair of species $(a,b)$, the balance functions
combine the four pair folders $(e,a,b)$, $(e,a,\bar b)$, $(e,\bar a,b)$, $(e,\bar a,\bar b)$ and the four single
folders $(e,a)$, $(e,\bar a)$, $(e,b)$, $(e,\bar b)$, and are written in a folder of their own,
\code{<e>\_<a>\_<b>\_BF}, next to the folders of the analyzer. Since a balance-function folder combines several
combinations of filters, its class (\code{ParticlePairBFHistos}, \code{ParticlePair3DBFHistos}) derives directly from
\code{HistogramFolder}, not from the pair folders.

\paragraph{One template for all phases.}
The tasks of the three phases derive from one template, \code{AnalysisHistos<HISTOS, RANK, SINGLES, BASE>}, which
manages a list of histogram folders of type \code{HISTOS}, one per combination of an event filter and \code{RANK}
particle filters, optionally with a list of single-particle folders of type \code{SINGLES} (for the pair analyses).
The template configures, creates or imports, scales, resets, saves and loads the folders, so that an analyzer
implements only the filling of an event and a calculator only the computation of its histograms. The base class
\code{BASE} is \code{EventProcessor} for the analyzers, through the alias
\code{EventProcessorHistos<HISTOS, RANK, SINGLES>}, and \code{AnalysisTask} for the calculators, through
\code{CalculatorHistos<HISTOS, RANK, SINGLES>}, which adds what all calculators share: the key \code{SOURCE\_TASK},
naming the analyzer whose folders they complete, and folders read from the input file of the job rather than created.
The calculators are not classes of their own but instances of three templates: \code{DerivedCalculator<HISTOS, RANK>},
for the analyses whose folders are each computed from their own base histograms (single-particle, global,
$\nudyn$, $\pT$--$\pT$); \code{PairDerivedCalculator<PAIR>}, for the pair analyses, whose pair folders also need
the single-particle folders; and \code{PairBFCalculator<PAIR, BF>}. Each is registered under the name of its task;
for instance,
{\footnotesize
\begin{verbatim}
GlobalDerivedCalculator
  = DerivedCalculator<GlobalHistos,1>
ParticlePairDerivedCalculator
  = PairDerivedCalculator<ParticlePairHistos>
ParticlePair3DBFCalculator
  = PairBFCalculator<ParticlePair3DHistos,
                     ParticlePair3DBFHistos>
\end{verbatim}}
\noindent The analyses thus differ only by their histograms, not by how they are called. RunCAP.py (Section~\ref{sec:run}) shows
the three phases as the stages of a workflow; it makes the derived-histogram stage from the analyzers of the first
stage, with the same filters and settings, and runs the stages as successive jobs.

% ---------------------------------------------------------------------------------------------------------------
\subsection{Histogram folders}
\label{sec:folders-analysis}

The histograms of an analysis are grouped in histogram folders, objects of classes derived from
\code{HistogramFolder}, itself a \code{Folder} (Section~\ref{sec:objects}) that owns its histograms;
Fig.~\ref{fig:hierarchy}(c) shows these classes. A folder creates its histograms from the histogram settings of its
task (binnings, ranges, and which histograms to fill) and keeps pointers to them for fast filling; it can also import
them from a file, scale them and reset them. Histograms have short names, \code{n1\_pt} or \code{n2\_DetaDphi} for
instance, since the folders tell the combinations apart. In the output file of a job, the folders of a task are the
sub-folders of a folder named after the task, so that a histogram is reached by a path such as
\code{PairAnalysis/ALL\_PiP\_PiM/R2\_DetaDphi}, both by programs and in the data browser of CapGui. Since folders of the same structure can be added
(Section~\ref{sec:objects}), the outputs of several jobs are added folder by folder.

% ---------------------------------------------------------------------------------------------------------------
\subsection{Merging jobs and sub-samples}
\label{sec:subsample}

Large productions are split into many jobs, run independently, for instance on a computing cluster, each writing its
own output file. Since the histograms of each job are normalised per accepted event, the jobs are combined by averaging
them, weighted by their numbers of events. The task \code{SubSampleStatCalculator} combines a set of jobs, one per
sub-directory of a path whose name matches a pattern, the jobs being treated as independent sub-samples. The files may
hold the histograms of any phase (filled, derived or balance functions), matched by their paths in the folders. With
$N$ sub-samples of values $x_i$ and weights $w_i$ ($W=\sum_i w_i$), bin by bin, it offers two estimates of the
statistical uncertainties, which matter for the quantities, such as the derived histograms, whose uncertainties are not
simply propagated from those of the bins:
\begin{itemize}
\item the sub-sample average $m=\sum_i w_i x_i/W$, with the uncertainty of the mean
  $\sigma^2 = \frac{N}{N-1}\sum_i w_i^2 (x_i-m)^2/W^2$;
\item the delete-one jackknife. Replica $i$ is the sample without sub-sample $i$. On the results of the jobs, its value
  is the average of the other sub-samples, $\theta_i=(Wm-w_ix_i)/(W-w_i)$, and $\theta=m$. Through the phases, the
  later phases of the workflow (derived histograms, balance functions) are run on the whole sample, the average of the
  filled histograms of the sub-samples, which gives $\theta$, and on each replica, which gives $\theta_i$: nonlinear
  observables such as correlation functions, ratios and balance functions are thus computed from (nearly) all the
  statistics, rather than averaged over sub-samples of small statistics. In both cases
  $\sigma^2=\frac{N-1}{N}\sum_i(\theta_i-\bar\theta)^2$, with $\bar\theta$ the mean of the $\theta_i$, and the
  bias-corrected values $N\theta-(N-1)\bar\theta$ are written to a file of their own.
\end{itemize}
The counts of events and of executed tasks of the jobs are summed and written with the result.

RunCAP.py (Section~\ref{sec:run}) runs a workflow as sub-samples and combines them. Each sub-sample has its own seeds:
in its workflow file, the seed parameters of the tasks (\code{Seed}, \code{SeedValue}) are increased by its number, and
so is the seed of the job (\code{JOB:SEED}). The sub-samples run on the computer, several at once, or on a Slurm
cluster, as a job array submitted over \code{ssh}, whose outputs are copied back and combined; jobs run in any other way
are combined from their output directories.
